\documentclass[10pt,a4paper]{amsart}
\usepackage{amsmath,amssymb,amsthm,graphicx}
\usepackage{verbatim}
\newenvironment{reference}{\comment}{\endcomment}
%\newenvironment{reference}{}{}
\newenvironment{slogan}{\comment}{\endcomment}
\newenvironment{history}{\comment}{\endcomment}

% For commutative diagrams we use Xy-pic
\usepackage[all]{xy}

% We use 2cell for 2-commutative diagrams.
\xyoption{2cell}
% Theorem environments.
%
\theoremstyle{plain}
\newtheorem{theorem}[subsection]{Theorem}
\newtheorem{proposition}[subsection]{Proposition}
\newtheorem{lemma}[subsection]{Lemma}

\theoremstyle{definition}
\newtheorem{definition}[subsection]{Definition}
\newtheorem{example}[subsection]{Example}
\newtheorem{exercise}[subsection]{Exercise}
\newtheorem{situation}[subsection]{Situation}

\theoremstyle{remark}
\newtheorem{remark}[subsection]{Remark}
\newtheorem{remarks}[subsection]{Remarks}

\numberwithin{equation}{subsection}

% Macros
%
\def\lim{\mathop{\mathrm{lim}}\nolimits}
\def\colim{\mathop{\mathrm{colim}}\nolimits}
\def\Spec{\mathop{\mathrm{Spec}}}
\def\Hom{\mathop{\mathrm{Hom}}\nolimits}
\def\Ext{\mathop{\mathrm{Ext}}\nolimits}
\def\SheafHom{\mathop{\mathcal{H}\!\mathit{om}}\nolimits}
\def\SheafExt{\mathop{\mathcal{E}\!\mathit{xt}}\nolimits}
\def\Sch{\mathit{Sch}}
\def\Mor{\mathop{\mathrm{Mor}}\nolimits}
\def\Ob{\mathop{\mathrm{Ob}}\nolimits}
\def\Sh{\mathop{\mathit{Sh}}\nolimits}
\def\NL{\mathop{N\!L}\nolimits}
\def\CH{\mathop{\mathrm{CH}}\nolimits}
\def\proetale{{pro\text{-}\acute{e}tale}}
\def\etale{{\acute{e}tale}}
\def\QCoh{\mathit{QCoh}}
\def\Ker{\mathop{\mathrm{Ker}}}
\def\Im{\mathop{\mathrm{Im}}}
\def\Coker{\mathop{\mathrm{Coker}}}
\def\Coim{\mathop{\mathrm{Coim}}}

% Boxtimes
%
\DeclareMathSymbol{\boxtimes}{\mathbin}{AMSa}{"02}

%
% Macros for moduli stacks/spaces
%
\def\QCohstack{\mathcal{QC}\!\mathit{oh}}
\def\Cohstack{\mathcal{C}\!\mathit{oh}}
\def\Spacesstack{\mathcal{S}\!\mathit{paces}}
\def\Quotfunctor{\mathrm{Quot}}
\def\Hilbfunctor{\mathrm{Hilb}}
\def\Curvesstack{\mathcal{C}\!\mathit{urves}}
\def\Polarizedstack{\mathcal{P}\!\mathit{olarized}}
\def\Complexesstack{\mathcal{C}\!\mathit{omplexes}}
% \Pic is the operator that assigns to X its picard group, usage \Pic(X)
% \Picardstack_{X/B} denotes the Picard stack of X over B
% \Picardfunctor_{X/B} denotes the Picard functor of X over B
\def\Pic{\mathop{\mathrm{Pic}}\nolimits}
\def\Picardstack{\mathcal{P}\!\mathit{ic}}
\def\Picardfunctor{\mathrm{Pic}}
\def\Deformationcategory{\mathcal{D}\!\mathit{ef}}
\usepackage{cleveref}
\crefname{originallemma}{original lemma}{original lemmas}
\Crefname{originallemma}{Original lemma}{Original lemmas}
\makeatletter
\newlabel{lemma-locally-constant-finite-type}{{093N}{source}{}{sourcepointer-1}{}}
\newlabel{lemma-locally-constant-finite-type@cref}{{[originallemma][0][]093N}{[1][0][]source}{}{}{}}
\makeatother
\title{\$\textbackslash{}SheafHom\$ and \$\textbackslash{}mathit\{Isom\}\$ of locally constant sheaves of modules are locally constant}
\author{The Stacks Project Authors}
\date{stacks-project:sites-modules.tex:local-snapshot-20261012; GNU FDL 1.2 or later}
\begin{document}
\maketitle
\noindent\textit{Editorial note.} All delivered mathematical text is the author's own, copied verbatim from the pinned Stacks Project chapter file sites-modules.tex; no statement or proof step was rewritten, reordered or removed.\par\medskip
\noindent\textit{Editorial note.} Span recomputation by label pairing (LANE\_BRIEF 10.2): in this chapter the author's proof environment follows the statement environment as a standalone block; the spec covers the complete environment pair (statement span, proof span) as recorded in the pool candidate and re-verified against the pinned source (sm-twins-20261014.json span sanity: all OK).\par\medskip
\noindent\textit{Editorial note.} Occupation scan: combined claim\_key + statement\_anchor + generic-line scan over all registration files (active and withdrawn) found hard occupation only from the 47 pre-existing lane-s11 S11-sites-mod-* cards (same pinned snapshot, delivered) and cross-chapter citation noise, 0 other-lane hits for sites-modules.tex (sm-occupied.json, 2026-10-14); the RESUME-pool dispatch unlocks the remaining face.\par\medskip
\ifdefined\hypertarget\hypertarget{sourcepointer-1}{}\fi
\noindent\textit{Source reference.} Original lemma 093N: sources/stacks-chow/tags.txt, source lines 11521--11521. Cross-chapter standard prerequisite cited by the author inside the counted span; it is used as a pointer and its own proof is not part of this item.\par\medskip
\begin{lemma}
\label{lemma-locally-constant}
Let $\mathcal{C}$ be a site. Let $\Lambda$ be a ring.
Let $M$, $N$ be $\Lambda$-modules.
Let $\mathcal{F}, \mathcal{G}$ be a locally constant sheaves of
$\Lambda$-modules.
\begin{enumerate}
\item If $M$ is of finite presentation, then
$$
\underline{\Hom_\Lambda(M, N)} =
\SheafHom_{\underline{\Lambda}}(\underline{M}, \underline{N})
$$
\item If $M$ and $N$ are both of finite presentation, then
$$
\underline{\text{Isom}_\Lambda(M, N)} =
\mathit{Isom}_{\underline{\Lambda}}(\underline{M}, \underline{N})
$$
\item If $\mathcal{F}$ is of finite presentation, then
$\SheafHom_{\underline{\Lambda}}(\mathcal{F}, \mathcal{G})$
is a locally constant sheaf of $\Lambda$-modules.
\item If $\mathcal{F}$ and $\mathcal{G}$ are both of finite presentation, then
$\mathit{Isom}_{\underline{\Lambda}}(\mathcal{F}, \mathcal{G})$
is a locally constant sheaf of sets.
\end{enumerate}
\end{lemma}

\begin{proof}
Proof of (1). Set $E = \Hom_\Lambda(M, N)$. We want to show the canonical map
$$
\underline{E}
\longrightarrow
\SheafHom_{\underline{\Lambda}}(\underline{M}, \underline{N})
$$
is an isomorphism. The module $M$ has a presentation
$\Lambda^{\oplus s} \to \Lambda^{\oplus t} \to M \to 0$.
Then $E$ sits in an exact sequence
$$
0 \to E \to \Hom_\Lambda(\Lambda^{\oplus t}, N) \to
\Hom_\Lambda(\Lambda^{\oplus s}, N)
$$
and we have similarly
$$
0 \to
\SheafHom_{\underline{\Lambda}}(\underline{M}, \underline{N}) \to
\SheafHom_{\underline{\Lambda}}(\underline{\Lambda^{\oplus t}}, \underline{N})
\to
\SheafHom_{\underline{\Lambda}}(\underline{\Lambda^{\oplus s}}, \underline{N})
$$
This reduces the question to the case where $M$ is a finite free module
where the result is clear.

\medskip\noindent
Proof of (3). The question is local on $\mathcal{C}$, hence we may assume
$\mathcal{F} = \underline{M}$ and $\mathcal{G} = \underline{N}$
for some $\Lambda$-modules $M$ and $N$.
By Lemma \ref{lemma-locally-constant-finite-type}
the module $M$ is of finite presentation. Thus the result follows from (1).

\medskip\noindent
Parts (2) and (4) follow from parts (1) and (3) and the
fact that $\mathit{Isom}$ can be viewed as the subsheaf of sections of
$\SheafHom_{\underline{\Lambda}}(\mathcal{F}, \mathcal{G})$
which have an inverse in
$\SheafHom_{\underline{\Lambda}}(\mathcal{G}, \mathcal{F})$.
\end{proof}

\end{document}
